\begin{abstract}
This paper examines whether trade agreement depth moderated the effect of the 2018 US steel and aluminum tariffs on import prices and values. Using a dynamic event study regression, I study how the percentage change of prices and values due to tariffs changes over a 12 month period following tariffs' imposition between countries with and without deep trade agreements, respectively. I measure trade agreement depth with the count of relevant policy areas to which countries' trade agreements with the United States apply. The findings explain whether expansive and institutionalized trade relationships increase resilience to unilateral tariffs.

% CLASS
This paper studies whether countries with deeper trade agreements – those covering more policy areas – saw a smaller impact on their trade with the United States as a result of the tariffs it imposed in 2018. Using monthly trade data, I track how import prices and volumes responded to tariffs over twelve months following their imposition, comparing that response across countries with deeper versus shallower trade relationships with the US. I hypothesize that countries with deep agreements prove more resilient because deep trade agreements reduce non-tariff barriers and lower overall bilateral trade costs; exporters operating under such agreements have a greater margin to absorb tariff-induced cost increases, sustaining trade volumes relative to countries without such institutional relationships. The results contribute to a growing literature on both how trade agreement architecture influences trade flows and on the dynamics of protectionism, offering evidence on whether having deep, institutionalized relationships with trading partners makes countries less vulnerable to sudden, unilateral tariffs.

% SUBMITTED TO CARROLL ROUND %
This paper examines whether the depth of trade agreements moderated the impact of the 2018 US steel and aluminum tariffs on import prices and volumes. While Amiti, Redding, and Weinstein (2020) establish that these tariffs were almost entirely borne by US importers with substantial reductions in trade volumes, whether this response varied with the structure of bilateral trade relationships remains an open question. I extend their dynamic event-study framework by introducing a triple interaction between event-time, the log change in tariff rates, and a binary indicator for trade agreement depth, yielding separate tariff elasticity paths for deep and non-deep agreement countries over the 18 months following tariff imposition. I hypothesize that deep-agreement countries prove more resilient, because deep trade agreements reduce non-tariff barriers and lower overall bilateral trade costs; exporters operating under such agreements have a greater margin to absorb tariff-induced cost increases, sustaining trade volumes relative to countries without such institutional relationships. Agreement depth is measured using the World Bank Deep Trade Agreements database, with a country classified as deep if its agreement with the United States contains ten or more legally enforceable horizontal policy provisions. The sample is restricted to steel and aluminum tariffs, which applied broadly across trading partners, avoiding the confounding influence of later tariff waves applied exclusively to China. I source trade flow and tariff data from US Census import data and US International Trade Commission tariff schedules. I estimate the model using Poisson Pseudo-Maximum Likelihood with high-dimensional fixed effects across multiple levels of product aggregation. The results contribute to a growing literature on both how trade agreement architecture influences trade flows and on the dynamics of protectionism, offering evidence on whether having deep, institutionalized relationships with trading partners makes countries less vulnerable to sudden, unilateral tariffs.

% SUBMITTED TO CCAS %
This paper studies whether countries with deeper trade agreements – those covering more policy areas – saw a smaller impact on their trade with the United States as a result of the tariffs it imposed in 2018. Using monthly trade data, I track how import prices and volumes responded to tariffs over twelve months following their imposition, comparing that response across countries with deeper versus shallower trade relationships with the US. The results offer evidence on whether having deep, institutionalized relationships with trading partners makes countries less vulnerable to sudden, unilateral tariffs.

help understand whether more comprehensive and institutionalized trade relationships enhance resilience to the disruptions caused by sudden, unilateral tariffs.

help understand whether deeper ties and institutionalized trade relations enhance resilience to the disruptions caused by sudden, unilateral tariffs.

The results speak to a practical question in international trade policy: Does having a deep, institutionalized relationship with a trading partner make a country less vulnerable when that partner turns protectionist?

Using monthly US trade data, I track how import prices and volumes responded to the tariffs over the twelve months following their imposition, and compare that response across countries with deeper versus shallower trade relationships with the US. I measure the depth of each country's trade agreement by counting the number of distinct policy areas it covers. The results shed light on whether strong, institutionalized trade relationships offer a buffer against the disruptions caused by sudden, unilateral tariff increases.

 using horizontal depth indicators from the World Bank Deep Trade Agreements database and use monthly US customs data at the HTS10 product level to measure import prices and values. The findings help illuminate how institutionalized trade relationships shape responses to unilaterally imposed tariffs.

Intro: context, problem, motivation 

Hypothesis: thesis, intentions 

Method: design, procedures, assumptions, data 

Results 

Conclusion: interpretation, inferences, implications, applications 

No: citations (unless doesn't make sense without mentioning it), details of findings
\end{abstract}
